\documentclass[a4paper,11pt]{article}
\usepackage[utf8]{inputenc}
\usepackage[spanish]{babel}
\usepackage{geometry}
\usepackage{enumitem}
\usepackage{sectsty}
\usepackage{titlesec}
\usepackage{tikz}
\usepackage{array}

\geometry{
    a4paper,
    left=20mm,
    top=18mm,
    right=20mm,
    bottom=20mm
}

% Sin números de página
\pagestyle{empty}

% Estilo de secciones con números en círculos
\newcommand{\circled}[1]{%
    \tikz[baseline=(char.base)]{%
        \node[shape=circle,draw,inner sep=0.5pt,fill=black,text=white,font=\bfseries\footnotesize] (char) {#1};%
    }%
}

% Formato de secciones
\titleformat{\section}
{\normalfont\Large\bfseries}
{\circled{\thesection}}{0.5em}{}

\titlespacing*{\section}{0pt}{12pt}{8pt}

\titleformat{\subsection}
{\normalfont\normalsize\bfseries}
{}{0em}{}

\titlespacing*{\subsection}{8pt}{4pt}{4pt}

\begin{document}

% Encabezado del curso
\begin{center}
\small{Universidad de Granada}

\vspace{0.3cm}

\textbf{\large{Producción e Interacción Oral - Nivel 7}}

\vspace{0.4cm}

\small
\begin{tabular}{@{}ll@{}}
\textbf{Grupos:} & A (Aula 24) y B (Aula 08) \\
\textbf{Centro:} & Centro de Lenguas Modernas / Edificio A \\
\textbf{Profesor:} & Javier Benítez Lainez \\
\textbf{Tutorías:} & Martes 9:00-10:00 y 14:30-15:30 \\
& Sala de profesores / Email: benitezl@go.ugr.es
\end{tabular}
\end{center}

\vspace{0.5cm}
\section*{Rúbrica de Autoevaluación de Debate}

\vspace{0.4cm}

\section{Instrumento de Reflexión y Auto-crítica C1}

\vspace{0.2cm}

---

\vspace{0.2cm}

\section{INTRODUCCIÓN}

\vspace{0.2cm}

\subsection{Propósito de esta rúbrica}

Esta rúbrica está diseñada para ayudarte a evaluar TU PROPIO desempeño en debates de forma honesta y constructiva. La autoevaluación es una herramienta esencial de aprendizaje en el nivel C1.

\vspace{0.2cm}

\subsection{Importancia de la autoevaluación}

\begin{itemize}

  \item \textbf{Metacognición}: Desarrolla conciencia sobre tu propio aprendizaje

  \item \textbf{Identificación de áreas de mejora}: Detecta debilidades específicas

  \item \textbf{Seguimiento de progreso}: Mide tu evolución a lo largo del curso

  \item \textbf{Autocrítica constructiva}: Aprende de tus propios errores

  \item \textbf{Independencia}: Desarrolla capacidad de auto-corrección

\end{itemize}

\vspace{0.2cm}

\subsection{Cómo usar esta rúbrica}

1. \textbf{Inmediatamente después del debate}: Reflexiona mientras está fresco

2. \textbf{Si es posible, revisa la grabación}: Objetividad adicional

3. \textbf{Sé honesto/a}: No seas ni demasiado crítico ni demasiado indulgente

4. \textbf{Específico/a}: Identifica ejemplos concretos

5. \textbf{Constructivo/a}: Enfócate en mejora, no solo en crítica

\vspace{0.2cm}

---

\vspace{0.2cm}

\section{PARTE 1: REFLEXIÓN INMEDIATA}

\vspace{0.2cm}

\subsection{Preguntas de reflexión post-debate}

\vspace{0.2cm}

\textbf{1. ¿Cómo te sentiste durante el debate?}

\begin{itemize}

  \item ☐ Muy confiada/o

  \item ☐ Algo confiada/o

  \item ☐ Nerviosa/o pero控制ada/o

  \item ☐ Muy nerviosa/o

  \item ☐ Aterrorizada/o

\end{itemize}

\vspace{0.2cm}

\textbf{Comentarios:} _________________________

\vspace{0.2cm}

\textbf{2. ¿Qué parte del debate sentiste que fue tu MEJOR momento?}

\begin{itemize}

  \item _________________________

\end{itemize}

\vspace{0.2cm}

\textbf{3. ¿Qué parte del辩论 sentiste que fue tu PEOR momento?}

\begin{itemize}

  \item _________________________

\end{itemize}

\vspace{0.2cm}

\textbf{4. ¿Cuál fue el momento más desafiante? ¿Por qué?}

\begin{itemize}

  \item _________________________

\end{itemize}

\vspace{0.2cm}

\textbf{5. ¿Hubo algo inesperado que ocurriera? ¿Cómo lo manejaste?}

\begin{itemize}

  \item _________________________

\end{itemize}

\vspace{0.2cm}

---

\vspace{0.2cm}

\section{PARTE 2: AUTOEVALUACIÓN DETALLADA}

\vspace{0.2cm}

\subsection{Sección 2.1: Preparación}

\vspace{0.2cm}

| Aspecto | Puntuación (1-4) | Evidencia/Ejemplo | Área de mejora |

|---------|------------------|-------------------|----------------|

| \textbf{Investigación del tema} | | | |

| | 4=Investigación profunda<br>3=Investigación buena<br>2=Investigación básica<br>1=Sin investigación | Ejemplo: _________________________ | _________________________ |

| \textbf{Preparación de argumentos} | | | |

| | 4=3+ argumentos bien preparados<br>3=2-3 argumentos<br>2=1-2 argumentos<br>1=Sin preparación | Ejemplo: _________________________ | _________________________ |

| \textbf{Anticipación de contraargumentos} | | | |

| | 4=Anticipé múltiples contraargumentos<br>3=Anticipé algunos<br>2=Anticipé pocos<br>1=No anticipé nada | Ejemplo: _________________________ | _________________________ |

| \textbf{Organización de notas} | | | |

| | 4=Notas bien organizadas y fáciles de usar<br>3=Notas organizadas<br>2=Notas algo desorganizadas<br>1=Sin notas o muy caóticas | Ejemplo: _________________________ | _________________________ |

\vspace{0.2cm}

\textbf{Puntuación total preparación: _____ / 16}

\vspace{0.2cm}

\textbf{Reflexión sobre preparación:}

\begin{itemize}

  \item ¿Qué funcionó bien en tu preparación?

\end{itemize}

_________________________

\begin{itemize}

  \item ¿Qué mejorarías la próxima vez?

\end{itemize}

_________________________

\vspace{0.2cm}

---

\vspace{0.2cm}

\subsection{Sección 2.2: Argumentación}

\vspace{0.2cm}

| Aspecto | Puntuación (1-4) | Evidencia/Ejemplo | Área de mejora |

|---------|------------------|-------------------|----------------|

| \textbf{Calidad de argumentos} | | | |

| | 4=Argumentos sofisticados y profundos<br>3=Argumentos sólidos<br>2=Argumentos básicos<br>1=Argumentos pobres | Ejemplo específico de tu MEJOR argumento: _________________________ | _________________________ |

| \textbf{Evidencia presentada} | | | |

| | 4=Evidencia específica y variada<br>3=Evidencia presente<br>2=Evidencia escasa<br>1=Sin evidencia | Ejemplo: _________________________ | _________________________ |

| \textbf{Refutación de oponente} | | | |

| | 4=Refutaciones excelentes<br>3=Refutaciones buenas<br>2=Refutaciones básicas<br>1=Sin refutaciones efectivas | Ejemplo de una refutación que hiciste: _________________________ | _________________________ |

| \textbf{Variedad de tipos de argumentos} | | | |

| | 4=Usé 4+ tipos (autoridad, causa-efecto, etc.)<br>3=Usé 3 tipos<br>2=Usé 1-2 tipos<br>1=Usé solo 1 tipo | Tipos que usaste: _________________________ | _________________________ |

\vspace{0.2cm}

\textbf{Puntuación total argumentación: _____ / 16}

\vspace{0.2cm}

\textbf{Reflexión sobre argumentación:}

\begin{itemize}

  \item ¿Qué argumento tuyo fue más efectivo? ¿Por qué?

\end{itemize}

_________________________

\begin{itemize}

  \item ¿Qué argumento podría haber mejorado? ¿Cómo?

\end{itemize}

_________________________

\vspace{0.2cm}

---

\vspace{0.2cm}

\subsection{Sección 2.3: Expresión Oral}

\vspace{0.2cm}

| Aspecto | Puntuación (1-4) | Evidencia/Ejemplo | Área de mejora |

|---------|------------------|-------------------|----------------|

| \textbf{Fluidez} | | | |

| | 4=Muy fluida, natural<br>3=Generalmente fluida<br>2=Algunas pausas<br>1=Muy interrumpida | Ejemplo: ¿Te trabaste en algún momento? ¿Dónde? _________________________ | _________________________ |

| \textbf{Vocabulario} | | | |

| | 4=Preciso, variado, nivel C1+<br>3=Adecuado, nivel C1<br>2=Básico, nivel B2+<br>1=Limitado, nivel B2- | Palabras o expresiones que usaste bien: _________________________ | _________________________ |

| \textbf{Gramática} | | | |

| | 4=Excelente control, errores mínimos<br>3=Buen control, algunos errores<br>2=Control aceptable, errores frecuentes<br>1=Pobre control | Errores gramaticales que notaste: _________________________ | _________________________ |

| \textbf{Pronunciación/Entonación} | | | |

| | 4=Clara, natural, énfasis apropiado<br>3=Clara, generally apropiada<br>2=Comprensible pero con effort<br>3=Difícil de entender | ¿Cómo sonó tu voz? ¿Confiable? ¿Nerviosa? _________________________ | _________________________ |

\vspace{0.2cm}

\textbf{Puntuación total expresión: _____ / 16}

\vspace{0.2cm}

\textbf{Reflexión sobre expresión:}

\begin{itemize}

  \item ¿En qué aspecto lingüístico te sentiste más segura/o?

\end{itemize}

_________________________

\begin{itemize}

  \item ¿Qué aspecto lingüístico necesita más trabajo?

\end{itemize}

_________________________

\vspace{0.2cm}

---

\vspace{0.2cm}

\subsection{Sección 2.4: Organización}

\vspace{0.2cm}

| Aspecto | Puntuación (1-4) | Evidencia/Ejemplo | Área de mejora |

|---------|------------------|-------------------|----------------|

| \textbf{Estructura (apertura, desarrollo, refutación, cierre)} | | | |

| | 4=Estructura clara y equilibrada<br>3=Estructura presente<br>2=Estructura básica<br>1=Sin estructura clara | ¿Cuál sección fue tu mejor? ¿Cuál peor? _________________________ | _________________________ |

| \textbf{Uso de conectores} | | | |

| | 4=Variados y apropiados<br>3=Adecuados<br>2=Básicos<br>1=Escasos o incorrectos | Conectores que usaste bien: _________________________ | _________________________ |

| \textbf{Gestión del tiempo} | | | |

| | 4=Perfectamente ajustada/o<br>3=Generalmente ajustada/o<br>2=Algo desajustada/o<br>1=Muy desajustada/o | ¿Te quedaste corta/o o te excediste? ¿En qué parte? _________________________ | _________________________ |

| \textbf{Coherencia global} | | | |

| | 4=Muy coherente y unificado<br>3=Coherente<br>2=Parcialmente coherente<br>1=Poco coherente | ¿Todas tus ideas conectaron bien? _________________________ | _________________________ |

\vspace{0.2cm}

\textbf{Puntuación total organización: _____ / 16}

\vspace{0.2cm}

\textbf{Reflexión sobre organización:}

\begin{itemize}

  \item ¿Qué aspecto organizacional funcionó mejor?

\end{itemize}

_________________________

\begin{itemize}

  \item ¿Cómo mejorarías tu organización la próxima vez?

\end{itemize}

_________________________

\vspace{0.2cm}

---

\vspace{0.2cm}

\subsection{Sección 2.5: Estrategias Comunicativas}

\vspace{0.2cm}

| Aspecto | Puntuación (1-4) | Evidencia/Ejemplo | Área de mejora |

|---------|------------------|-------------------|----------------|

| \textbf{Interacción con oponente} | | | |

| | 4=Escucha activa, respeta, responde bien<br>3=Generalmente interactúa bien<br>2=Interacción limitada<br>1=Poca o ninguna interacción | ¿Escuchaste realmente a tu oponente? _________________________ | _________________________ |

| \textbf{Lenguaje no verbal} | | | |

| | 4=Contacto visual, gestos, postura excelentes<br>3=Generalmente buenos<br>2=Básicos<br>1=Pobres o nerviosos | ¿Qué notaste sobre tu lenguaje corporal? _________________________ | _________________________ |

| \textbf{Recursos retóricos} | | | |

| | 4=Varios recursos bien usados<br>3=Algunos recursos<br>2=Pocos recursos<br>1=Sin recursos | ¿Usaste preguntas retóricas, análogías, tripletas? _________________________ | _________________________ |

| \textbf{Convicción y persuasión} | | | |

| | 4=Muy convincente y persuasiva/o<br>3=Algo convincente<br>2=Poco convincente<br>1=No convincente | ¿Te sentiste convencida/o? ¿Por qué sí o no? _________________________ | _________________________ |

\vspace{0.2cm}

\textbf{Puntuación total estrategias: _____ / 16}

\vspace{0.2cm}

\textbf{Reflexión sobre estrategias:}

\begin{itemize}

  \item ¿Qué estrategia comunicativa usaste mejor?

\end{itemize}

_________________________

\begin{itemize}

  \item ¿Qué estrategia comunicativa necesita más desarrollo?

\end{itemize}

_________________________

\vspace{0.2cm}

---

\vspace{0.2cm}

\section{PARTE 3: ANÁLISIS DE FORTALEZAS Y DEBILIDADES}

\vspace{0.2cm}

\subsection{Fortalezas principales (Top 3)}

\vspace{0.2cm}

\textbf{1. } _________________________

\begin{itemize}

  \item Evidencia: _________________________

\end{itemize}

\vspace{0.2cm}

\textbf{2. } _________________________

\begin{itemize}

  \item Evidencia: _________________________

\end{itemize}

\vspace{0.2cm}

\textbf{3. } _________________________

\begin{itemize}

  \item Evidencia: _________________________

\end{itemize}

\vspace{0.2cm}

---

\vspace{0.2cm}

\subsection{Áreas de mejora principales (Top 3)}

\vspace{0.2cm}

\textbf{1. } _________________________

\begin{itemize}

  \item Plan de acción: _________________________

\end{itemize}

\vspace{0.2cm}

\textbf{2. } _________________________

\begin{itemize}

  \item Plan de acción: _________________________

\end{itemize}

\vspace{0.2cm}

\textbf{3. } _________________________

\begin{itemize}

  \item Plan de acción: _________________________

\end{itemize}

\vspace{0.2cm}

---

\vspace{0.2cm}

\section{PARTE 4: COMPARACIÓN CON DEBATES ANTERIORES}

\vspace{0.2cm}

\subsection{Si has hecho debates anteriores, compara:}

\vspace{0.2cm}

| Aspecto | Debate anterior | Este debate | ¿Mejoró? (Sí/No/Parcial) |

|---------|----------------|-------------|--------------------------|

| Preparación | | | |

| Argumentación | | | |

| Fluidez | | | |

| Vocabulario | | | |

| Gramática | | | |

| Confianza | | | |

| Gestión del tiempo | | | |

\vspace{0.2cm}

\textbf{Reflexión sobre progreso:}

\begin{itemize}

  \item ¿En qué áreas has mejorado más?

\end{itemize}

_________________________

\begin{itemize}

  \item ¿Qué áreas siguen necesitando trabajo?

\end{itemize}

_________________________

\vspace{0.2cm}

---

\vspace{0.2cm}

\section{PARTE 5: OBJETIVOS PARA EL PRÓXIMO DEBATE}

\vspace{0.2cm}

\subsection{Específicos, medibles, alcanzables, relevantes, temporales (SMART)}

\vspace{0.2cm}

\textbf{Objetivo 1:}

\begin{itemize}

  \item \textbf{Qué mejorarás:} _________________________

  \item \textbf{Cómo lo harás:} _________________________

  \item \textbf{Cómo medirás el éxito:} _________________________

\end{itemize}

\vspace{0.2cm}

\textbf{Objetivo 2:}

\begin{itemize}

  \item \textbf{Qué mejorarás:} _________________________

  \item \textbf{Cómo lo harás:} _________________________

  \item \textbf{Cómo medirás el éxito:} _________________________

\end{itemize}

\vspace{0.2cm}

\textbf{Objetivo 3:}

\begin{itemize}

  \item \textbf{Qué mejorarás:} _________________________

  \item \textbf{Cómo lo harás:} _________________________

  \item \textbf{Cómo medirás el éxito:} _________________________

\end{itemize}

\vspace{0.2cm}

---

\vspace{0.2cm}

\section{PARTE 6: REFLEXIÓN FINAL}

\vspace{0.2cm}

\subsection{Pregunta de reflexión profunda}

\vspace{0.2cm}

\textbf{¿Qué aprendiste sobre TI MISMA/O como comunicadora/or durante este debate?}

\vspace{0.2cm}

Respuesta: _________________________

\vspace{0.2cm}

\textbf{¿Qué harías diferente si pudieras repetir este debate?}

\vspace{0.2cm}

Respuesta: _________________________

\vspace{0.2cm}

\textbf{¿Qué consejo le darías a un/a estudiante que va a hacer su primer debate?}

\vspace{0.2cm}

Respuesta: _________________________

\vspace{0.2cm}

---

\vspace{0.2cm}

\section{PUNTUACIÓN TOTAL DE AUTOEVALUACIÓN}

\vspace{0.2cm}

\textbf{Preparación:} _____ / 16

\textbf{Argumentación:} _____ / 16

\textbf{Expresión Oral:} _____ / 16

\textbf{Organización:} _____ / 16

\textbf{Estrategias Comunicativas:} _____ / 16

\vspace{0.2cm}

\textbf{TOTAL: _____ / 80}

\vspace{0.2cm}

\subsection{Conversión a auto-nota}

\vspace{0.2cm}

\begin{itemize}

  \item \textbf{72-80} (90-100%): \textbf{Autoevaluación Excelente}

  \item \textbf{56-71} (70-89%): \textbf{Autoevaluación Buena}

  \item \textbf{40-55} (50-69%): \textbf{Autoevaluación Suficiente}

  \item \textbf{0-39} (0-49%): \textbf{Autoevaluación Insuficiente}

\end{itemize}

\vspace{0.2cm}

\textbf{Tu auto-nota: _____}

\vspace{0.2cm}

---

\vspace{0.2cm}

\section{COMPARACIÓN CON NOTA DEL PROFESOR}

\vspace{0.2cm}

Después de recibir la nota del profesor/a:

\vspace{0.2cm}

\textbf{Nota profesor/a:} _____

\textbf{Tu auto-nota:} _____

\textbf{Diferencia:} +/- _____ puntos

\vspace{0.2cm}

\textbf{¿Por qué crees que hay diferencia?}

\begin{itemize}

  \item _________________________

\end{itemize}

\vspace{0.2cm}

\textbf{¿Te evaluaste demasiado duramente o demasiado indulgentemente?}

\begin{itemize}

  \item _________________________

\end{itemize}

\vspace{0.2cm}

---

\vspace{0.2cm}

\section{FIRMA}

\vspace{0.2cm}

\textbf{Estudiante:} _________________________

\textbf{Fecha del debate:} _________________________

\textbf{Tema del debate:} _________________________

\textbf{Fecha de autoevaluación:} _________________________

\vspace{0.2cm}

---

\vspace{0.2cm}

\textbf{Sesión 8: Autoevaluación de Debate}

\textbf{Nivel C1 - Español Oral Avanzado}

\vspace{0.2cm}

\end{document}
